\section{Semistable and unstable cohomology}\label{sec:asymptotics}

This section describes the long-time behaviour when $\mathcal H$ is not
polystable. For semistable $\mathcal H$ the curve of normalized metrics
stays at bounded distance from a trajectory of a gradient flow of the
Kempf--Ness function, whose asymptotics are known
(Theorem~\ref{thm:comparison}). With no hypothesis on $\mathcal H$, the mass
converges to the mass of the Harder--Narasimhan filtration, the filtration
of the spaces by the exponential rate of decay of the metric is that
filtration, and its factors, which are semistable, follow the semistable
theory (Theorem~\ref{thm:hn}).

Throughout the section $Q$ is a finite quiver without oriented cycles,
$\mathcal H$ is a nonzero representation of $Q$, with
$Z(\mathcal H)=|Z|e^{\ii\varphi}$, and $\mathfrak h(t)$,
$t\geq\tau_0$, is a curve of metrics on $\mathcal H$ as in
Proposition~\ref{prop:perturbed}: it is continuously differentiable, its
velocity is $\xi_i=G_i(\mu_i)-\delta_i$ with continuous
$\mathfrak h$-Hermitian $\delta_i$,
$\norm{\delta(t)}\leq c_\delta e^{-\kappa(t-\tau_0)}$ with $\kappa>0$, and the
sum of the squared norms of the arrows is at most $\bar\cN$. We call this
\emph{hypothesis} (P). No stability of $\mathcal H$ is assumed in (P). It
holds, with $\delta=0$, for the normalized phase flow
$\mathfrak h=e^{\varphi t}h$ on a representation of a quiver without oriented
cycles (Proposition~\ref{prop:reparrows}); and it holds for the normalized
metrics induced on the cohomology of a complex with a surjective or an
injective differential at one vertex, for every cohomology representation:
the proof of Theorem~\ref{thm:convergence}(i) and
Theorem~\ref{thm:decay} use no stability, and Remark~\ref{rem:vanishing}
covers the complexes excluded there. For an injective differential one
passes to the dual complex (Proposition~\ref{prop:injective}): the
normalized induced metrics of the two are dual to each other, so their
velocities are $\xi$ and $-\xi^{\mathsf T}$; the moment map changes sign
under duality; and
$\varphi-\arg(z_i+u)=-\bigl[(\pi-\varphi)-\arg(-\overline{z_i}-u)\bigr]$
for real $u$. So (P) for the dual gives (P).

We write $\Phi^{\mathcal H}_i=\arg(z_i+\mu_i)=\varphi-G_i(\mu_i)$ for the
phase of the representation $\mathcal H$ with the metric $\mathfrak h$, a
Hermitian matrix in an orthonormal frame; the function
$u\mapsto\arg(z_i+u)$ is decreasing, with derivative in $[-1/y_i,0)$. In
orthonormal frames of $\mathfrak h(t)$ whose coordinate maps $g_i$, with
$|g_iv|^2=\mathfrak h_i(v,v)$, satisfy $\dot g_ig_i^{-1}=\frac12\xi_i$, the
arrows evolve by
\begin{equation}\label{eq:betaflow}
 \dot\beta_a=\tfrac12\bigl(\beta_a(\Phi^{\mathcal H}_{\mathsf s(a)}+\delta_{\mathsf s(a)})
 -(\Phi^{\mathcal H}_{\mathsf t(a)}+\delta_{\mathsf t(a)})\beta_a\bigr),
\end{equation}
and the eigenvalues of all $\Phi^{\mathcal H}_i$ lie in a compact interval
$[\psi_-,\psi_+]\subset(0,\pi)$ determined by the charges and $\bar\cN$. Two
facts on functions of Hermitian matrices are used. If $g$ is a continuously
differentiable function on $\RR$ with $|g'|\leq c$, then
$\norm{g(A)-g(B)}\leq c\norm{A-B}$ for Hermitian $A,B$ and the
Hilbert--Schmidt norm; and if $g$ is nonincreasing, then
$\tr\bigl((g(A)-g(B))(A-B)\bigr)\leq0$. Both follow from the identities
$\inner{u_j}{(g(A)-g(B))v_l}=(g(a_j)-g(b_l))\inner{u_j}{v_l}$ and
$\inner{u_j}{(A-B)v_l}=(a_j-b_l)\inner{u_j}{v_l}$ for orthonormal
eigenvectors $u_j$ of $A$ and $v_l$ of $B$ with eigenvalues $a_j$, $b_l$.

\subsection{Semistable cohomology: comparison with the gradient flow}
Put
\[
 \mathsf m_i=\frac1{G_i'(\theta_i)}=\frac{y_i}{\sigma^2}\qquad(i\in Q_0),
\]
and let $d_{\mathsf m}$ be the distance of the Riemannian metric
$\norm\xi_{\mathsf m}^2=\sum_i\mathsf m_i\tr\xi_i^2$ on $\mathcal X$. With
this metric $\mathcal X$ is the product of the symmetric spaces of positive
definite Hermitian matrices, the $i$-th factor with its invariant metric
multiplied by $\mathsf m_i$; it is complete, simply connected and
nonpositively curved, its geodesics are again
$u\mapsto\mathfrak he^{u\xi}$, and
$d_{\mathsf m}(\mathfrak h,\mathfrak h')^2=\sum_i\mathsf m_i\norm{\log(\mathfrak h_i^{-1/2}\mathfrak h_i'\mathfrak h_i^{-1/2})}^2$.
By Lemma~\ref{lem:variation} the gradient of $\cK_{\mathcal H}$ for this
metric is the tangent vector $(-(\mu_i-\theta_i)/\mathsf m_i)_i$. The
\emph{weighted gradient flow} is
\begin{equation}\label{eq:gradflow}
 \mathfrak g_i^{-1}\dot{\mathfrak g}_i=\frac{\mu_i(\mathfrak g)-\theta_i}{\mathsf m_i}\qquad(i\in Q_0).
\end{equation}
Its solutions exist for all later times: the sum $q$ of the squared norms of
the arrows satisfies
$\dot q=-\sum_i\tr(\mu_i\xi_i)\leq\sum_in_i\theta_i^2/(4\mathsf m_i)$ along
\eqref{eq:gradflow}, so the velocity is bounded on bounded intervals of
time. The flow \eqref{eq:gradflow} is the gradient flow studied in
\cite{hkkp23}: it is \cite[(5.11)]{hkkp23}, with the parameters
$\mathsf m_i$ and $\theta_i$ there.

\begin{lemma}\label{lem:quadratic}
For all real $u$ and all $i$,
$\bigl|G_i(u)-(u-\theta_i)/\mathsf m_i\bigr|\leq\frac{3\sqrt3}{16\,y_i^2}(u-\theta_i)^2$.
Consequently
$\norm{G_i(M)-(M-\theta_i)/\mathsf m_i}\leq\frac{3\sqrt3}{16\,y_i^2}\norm{M-\theta_i}^2$
for every Hermitian matrix $M$.
\end{lemma}

\begin{proof}
$G_i''(u)=-2y_i(x_i+u)/((x_i+u)^2+y_i^2)^2$, and
$2y|w|/(w^2+y^2)^2$ has its maximum $3\sqrt3/(8y^2)$ at $w^2=y^2/3$. Use
Taylor's formula at $\theta_i$, where $G_i=0$ and $G_i'=1/\mathsf m_i$. For a
Hermitian matrix apply this to the eigenvalues $u_k$, and
$(\sum_k(u_k-\theta_i)^4)^{1/2}\leq\sum_k(u_k-\theta_i)^2$.
\end{proof}

\begin{lemma}\label{lem:hadamard}
Let $\mathcal M$ be a complete, simply connected Riemannian manifold of
nonpositive sectional curvature, with distance $d$, and $F$ a continuously
differentiable function on $\mathcal M$ that is convex along geodesics. Let
$y(t)$ be a solution of $\dot y=-\operatorname{grad}F(y)$ and $x(t)$ a
continuously differentiable curve with
$\dot x=-\operatorname{grad}F(x)+e(t)$ on an interval $[t_0,t_1]$. Then
$d(x(t),y(t))\leq d(x(t_0),y(t_0))+\int_{t_0}^t|e|$ for $t_0\leq t\leq t_1$.
\end{lemma}

\begin{proof}
On such a manifold any two points $p,p'$ are joined by a unique geodesic
$\gamma_{p,p'}:[0,1]\to\mathcal M$, which depends smoothly on $(p,p')$, and
$d(p,p')^2=\int_0^1|\gamma_{p,p'}'|^2$ is smooth. By the first variation
formula, with $\gamma_t=\gamma_{x(t),y(t)}$ and $\rho(t)=d(x(t),y(t))$,
\[
 \frac d{dt}\,\frac{\rho^2}2=\inner{\dot y}{\gamma_t'(1)}-\inner{\dot x}{\gamma_t'(0)}
 =-\inner{\operatorname{grad}F(y)}{\gamma_t'(1)}+\inner{\operatorname{grad}F(x)}{\gamma_t'(0)}
 -\inner e{\gamma_t'(0)} .
\]
The function $u\mapsto F(\gamma_t(u))$ is convex and differentiable, so its
derivative $\inner{\operatorname{grad}F(\gamma_t(u))}{\gamma_t'(u)}$ is
nondecreasing, and the right side is at most $|e|\,\rho$. For
$\epsilon>0$ this gives $\frac d{dt}\sqrt{\rho^2+\epsilon}\leq|e|$; integrate
and let $\epsilon\to0$.
\end{proof}

For $e=0$ the lemma says that the gradient flow does not increase distances;
Harada and Wilkin prove a related contraction estimate for the trace of the relative metric and its inverse \cite[Theorem~3.18]{hw11}; see also \cite[\S5.3]{hkkp23}.

\begin{theorem}[Comparison with the gradient flow]\label{thm:comparison}
Assume (P) and let $\mathcal H$ be semistable. Let $\mathfrak g(t)$,
$t\geq\tau_0$, be the solution of the weighted gradient flow
\eqref{eq:gradflow} with $\mathfrak g(\tau_0)=\mathfrak h(\tau_0)$. Then for
all $t\geq\tau_0$
\[
 d_{\mathsf m}\bigl(\mathfrak h(t),\mathfrak g(t)\bigr)\leq C_{\mathrm e}
 :=\max_i\frac{3\sqrt3}{16\,\sigma\,y_i^{3/2}}\Bigl(2\Gamma\bigl(\cK_{\mathcal H}(\mathfrak h(\tau_0))-\cK_*\bigr)
 +\frac{\Gamma^2c_\delta^2}{2\kappa}\Bigr)
 +\frac{c_\delta}\kappa\,\max_i\sqrt{\mathsf m_i},
\]
with $\Gamma$ as in \eqref{eq:Gamma}; and, as Hermitian forms,
$e^{-c}\mathfrak g_i(t)\leq\mathfrak h_i(t)\leq e^{c}\mathfrak g_i(t)$ with
$c=C_{\mathrm e}/\min_i\sqrt{\mathsf m_i}$.
\end{theorem}

\begin{proof}
The velocity of $\mathfrak h$ is minus the gradient of $\cK_{\mathcal H}$ for
the metric $\norm\cdot_{\mathsf m}$ plus the tangent vector
$e_i=[G_i(\mu_i)-(\mu_i-\theta_i)/\mathsf m_i]-\delta_i$. By
Lemma~\ref{lem:quadratic},
$\norm e_{\mathsf m}\leq\max_i\frac{3\sqrt3}{16\,\sigma y_i^{3/2}}\norm{\mu-\theta}^2+\max_i\sqrt{\mathsf m_i}\,\norm\delta$.
By Proposition~\ref{prop:perturbed}(i) and Young's inequality,
$\frac d{dt}\cK_{\mathcal H}(\mathfrak h(t))\leq-\frac1{2\Gamma}\norm{\mu-\theta}^2+\frac\Gamma2\norm\delta^2$,
and $\cK_{\mathcal H}\geq\cK_*>-\infty$ because $\mathcal H$ is semistable
(Lemma~\ref{lem:king}(a)); so
$\int_{\tau_0}^\infty\norm{\mu-\theta}^2\leq2\Gamma(\cK_{\mathcal H}(\mathfrak h(\tau_0))-\cK_*)+\Gamma^2c_\delta^2/(2\kappa)$.
Lemma~\ref{lem:hadamard} on $\mathcal X$ with the metric
$\norm\cdot_{\mathsf m}$ and $F=\cK_{\mathcal H}$, which is convex along the
geodesics $u\mapsto\mathfrak he^{u\xi}$ (Lemma~\ref{lem:variation}), gives
the first claim. The second follows from
$\norm{\log(\mathfrak g_i^{-1/2}\mathfrak h_i\mathfrak g_i^{-1/2})}\leq d_{\mathsf m}(\mathfrak h,\mathfrak g)/\sqrt{\mathsf m_i}$.
\end{proof}

\begin{corollary}\label{cor:presentations}
Two curves that satisfy (P) for the same semistable representation
$\mathcal H$ and the same charges remain at bounded distance from each
other. In particular this holds for the normalized metrics induced on the
cohomology of two complexes with a surjective or injective differential at
one vertex whose cohomology representations are identified by a fixed
isomorphism, and for any two initial metrics.
\end{corollary}

\begin{proof}
Each of the two curves is at bounded distance from a solution of
\eqref{eq:gradflow} by Theorem~\ref{thm:comparison}, and two solutions of
\eqref{eq:gradflow} do not move apart by Lemma~\ref{lem:hadamard} with $e=0$.
\end{proof}

The asymptotics of the flow \eqref{eq:gradflow} on a semistable
representation are the subject of \cite{hkkp23}. In the notation of
\cite[\S5.2]{hkkp23} the trace is $\tau(b)=\sum_i\mathsf m_i\tr b_i$, the
central element $\rho$ has the components $\theta_i/\mathsf m_i$, the flow is
\cite[(5.28)]{hkkp23}, and the standing assumption \cite[(5.59)]{hkkp23} is
the semistability of $\mathcal H$ in the sense of
Section~\ref{sec:convergence}. A curve $\mathfrak h$ is called there an
\emph{asymptotic solution} if $C^{-1}\mathfrak g\leq\mathfrak h\leq C\mathfrak g$
for large times, for some, equivalently for every, solution $\mathfrak g$
\cite[\S5.3]{hkkp23}. In \cite[\S5.4]{hkkp23} the lattice of
subrepresentations of $\mathcal H$ of the same phase carries the
\emph{iterated weight filtration}, a filtration labelled by finite sequences
$\lambda=(\lambda_1,\dots,\lambda_n)$ of real numbers and defined
algebraically from that lattice and the function
$\mathsf R\mapsto\sum_i\mathsf m_i\dim\mathsf R_i$
\cite[\S4 and (5.58)]{hkkp23}. After the arrows have been conjugated by a
suitable invertible tuple of matrices, with $1=\sum_\lambda p_\lambda$ the
orthogonal splitting given by the filtration, the function
\begin{equation}\label{eq:iterated}
 \mathfrak a(t)=\sum_{\lambda}t^{\lambda_1}(\log t)^{\lambda_2}\cdots\bigl(\log^{(n-1)}t\bigr)^{\lambda_n}\,p_\lambda
\end{equation}
is an asymptotic solution \cite[Corollary~5.10]{hkkp23}, and every solution
satisfies the bounds of \cite[Theorem~5.11]{hkkp23} on the steps of the
filtration.

\begin{corollary}[Iterated logarithms]\label{cor:iterated}
Assume (P) and let $\mathcal H$ be semistable. Then $\mathfrak h$ is an
asymptotic solution of the gradient flow \cite[(5.28)]{hkkp23} with the
weights $\mathsf m_i=y_i/\sigma^2$. Hence there are an invertible tuple of
matrices $b$ and a constant $C$ such that
$C^{-1}b^*\mathfrak a(t)b\leq\mathfrak h(t)\leq Cb^*\mathfrak a(t)b$ for large
$t$, where $\mathfrak a$ is the function \eqref{eq:iterated} for the arrows
of $\mathcal H$ conjugated by $b$; and the bounds of
\cite[Theorem~5.11]{hkkp23} hold for $\mathfrak h$. The weight function of
the iterated weight filtration is
\[
 \sum_i\mathsf m_i\dim\mathsf R_i=\frac{\Ima Z(\mathsf R)}{\sigma^2}=\frac{|Z(\mathsf R)|}\sigma
\]
on the subquotients $\mathsf R$ of $\mathcal H$ that are semistable of phase
$\varphi$: a constant multiple of the mass.
\end{corollary}

\begin{proof}
The first statement is the last statement of
Theorem~\ref{thm:comparison}; the others follow from the definition of an
asymptotic solution and the results of \cite{hkkp23} just quoted. For a
projection $p$ the inequalities
$e^{-c}\mathfrak g\leq\mathfrak h\leq e^{c}\mathfrak g$ give
$e^{-c}p\mathfrak gp\leq p\mathfrak hp\leq e^{c}p\mathfrak gp$.
\end{proof}

The first statement of the corollary can also be deduced from
\cite[Proposition~5.7]{hkkp23}, which says that a curve satisfying the
equation of the gradient flow up to an absolutely integrable term is an
asymptotic solution; the term is the vector $e$ of the proof of
Theorem~\ref{thm:comparison}. The bounded factor between $\mathfrak h$ and
\eqref{eq:iterated} is not shown to converge, also not in \cite{hkkp23} for
the gradient flow; it converges in Proposition~\ref{prop:weights} below and
in the example of \cite[Proposition~6.9]{wue26b}.
In the following example the exponents determine the ratio of the vertex
weights.

\begin{proposition}\label{prop:weights}
Let $Q$ be the quiver $0\to1$, let $z_0,z_1$ have the same argument
$\varphi$, and let $\mathcal H=(\CC\xrightarrow{b}\CC)$ with $b\neq0$; it is
semistable and not polystable, a nonsplit extension of the simple
representation at the vertex $0$ by the simple representation at the vertex
$1$. For every initial metric the normalized phase flow
$\mathfrak h=e^{\varphi t}h$ satisfies
\[
 \log\mathfrak h_0(t)=\frac{|z_1|}{|z_0|+|z_1|}\log t+c_0+o(1),\qquad
 \log\mathfrak h_1(t)=-\frac{|z_0|}{|z_0|+|z_1|}\log t+c_1+o(1)
\]
as $t\to\infty$, with constants $c_0,c_1$. The solutions of
\eqref{eq:gradflow} with arbitrary positive weights $\mathsf m_0',\mathsf m_1'$
have the exponents $\mathsf m_1'/(\mathsf m_0'+\mathsf m_1')$ and
$-\mathsf m_0'/(\mathsf m_0'+\mathsf m_1')$.
\end{proposition}

\begin{proof}
Here $\theta=0$ and $\mathsf m_i=|z_i|/\sigma$. The squared norm
$q=|b|^2\mathfrak h_1/\mathfrak h_0$ of the arrow gives $\mu_0=q$,
$\mu_1=-q$, so $(\log\mathfrak h_0)'=G_0(q)$, $(\log\mathfrak h_1)'=G_1(-q)$
and $(\log q)'=-(G_0(q)-G_1(-q))<0$. By Lemma~\ref{lem:quadratic},
$G_0(q)-G_1(-q)=\varkappa q+O(q^2)$ with
$\varkappa=1/\mathsf m_0+1/\mathsf m_1$. So $q$ decreases to zero,
$\frac d{dt}(1/q)=\varkappa+O(q)$, hence $1/q=\varkappa t+O(\log t)$ and
$q=1/(\varkappa t)+O(t^{-2}\log t)$. Then
$(\log\mathfrak h_0)'=q/\mathsf m_0+O(q^2)$ integrates to
$\frac1{\mathsf m_0\varkappa}\log t+c_0+o(1)$, and
$\frac1{\mathsf m_0\varkappa}=\frac{\mathsf m_1}{\mathsf m_0+\mathsf m_1}$;
likewise for $\mathfrak h_1$. For \eqref{eq:gradflow} with weights
$\mathsf m_i'$ the same computation applies with
$(\log\mathfrak g_0)'=q/\mathsf m_0'$, $(\log\mathfrak g_1)'=-q/\mathsf m_1'$.
\end{proof}

So on this representation the metric grows like $t$ to the power
$|Z(\mathsf S_1)|/|Z(\mathcal H)|$ on the quotient $\mathsf S_0$ and decays
like $t$ to the power $-|Z(\mathsf S_0)|/|Z(\mathcal H)|$ on the
subrepresentation $\mathsf S_1$, where $\mathsf S_i$ is the simple
representation at the vertex $i$; and a gradient flow whose weights are not
proportional to the masses $|z_i|$ is at logarithmically growing distance
from the phase flow. For a generalized Kronecker quiver with
$\alpha_0\neq\alpha_1$ the weights do not matter, because the stable
factors of a semistable representation then have proportional dimension
vectors.

\subsection{The mass and the Harder--Narasimhan filtration}
A representation $\mathsf R\neq0$ is semistable of phase
$\phi=\arg Z(\mathsf R)$, in the sense of Section~\ref{sec:convergence} with
the parameter $\theta$ formed with $\phi$, if and only if
$\arg Z(\mathsf R')\leq\phi$ for all nonzero subrepresentations
$\mathsf R'$. Every nonzero representation $\mathcal H$ has a unique
filtration
\[
 0=\mathcal H^0\subset\mathcal H^1\subset\dots\subset\mathcal H^{\mathsf n}=\mathcal H
\]
by subrepresentations whose factors
$\operatorname{gr}^k=\mathcal H^k/\mathcal H^{k-1}$ are semistable of phases
$\phi_1>\phi_2>\dots>\phi_{\mathsf n}$, the \emph{Harder--Narasimhan
filtration} for the stability condition of Section~\ref{sec:setting}
\cite[Proposition~2.4]{bri07}; a filtration by subrepresentations with
semistable factors of strictly decreasing phases is this one, because there
are no nonzero morphisms from a semistable representation to a semistable
representation of smaller phase. We call
$\mass_{\mathrm{HN}}(\mathcal H)=\sum_k|Z(\operatorname{gr}^k)|$ the
\emph{mass of the filtration}; it is the mass of the object $\mathcal H$ for
the stability condition. For a tuple of metrics on $\mathcal H$ let
$\mass(\mathfrak h)=\sum_i\tr|z_i+\mu_i|$ be the mass of the metrized
representation, $\Pi_k$ the orthogonal projections onto $\mathcal H^k$, and
\[
 \mathcal Y_k=\sum_a\norm{\Pi_{k,\mathsf t(a)}\beta_a(1-\Pi_{k,\mathsf s(a)})}^2
\]
the squared norm of the second fundamental form of $\mathcal H^k$.

\begin{proposition}[Dissipation of the mass]\label{prop:dissipation}
Assume (P). Then for $t\geq\tau_0$
\[
 \frac{d\,\mass(\mathfrak h(t))}{dt}\leq-\mathfrak D+4\bar\cN\norm\delta,
 \qquad
 \mathfrak D\geq\min(\sin\psi_-,\sin\psi_+)\sum_a\norm{\beta_a\Phi^{\mathcal H}_{\mathsf s(a)}-\Phi^{\mathcal H}_{\mathsf t(a)}\beta_a}^2,
\]
where
$\mathfrak D=\sum_a\sum_{j,l}|\inner{v_l}{\beta_au_j}|^2(p_j-q_l)(\cos q_l-\cos p_j)$
for orthonormal eigenvectors $u_j$ of $\Phi^{\mathcal H}_{\mathsf s(a)}$ and
$v_l$ of $\Phi^{\mathcal H}_{\mathsf t(a)}$ with eigenvalues $p_j$, $q_l$. In
particular the mass is nonincreasing along the phase flow of a
representation, for all charges.
\end{proposition}

\begin{proof}
The derivative of $u\mapsto|z_i+u|$ is $\cos\arg(z_i+u)$, so
$\frac d{dt}\mass=\sum_i\tr(\cos\Phi^{\mathcal H}_i\,\dot\mu_i)$. Insert
\eqref{eq:betaflow}; the result is linear in
$\Phi^{\mathcal H}+\delta$. The part with $\Phi^{\mathcal H}$ is, since
$\Phi^{\mathcal H}_i$ commutes with its cosine,
\[
 \sum_a\sum_{j,l}|\inner{v_l}{\beta_au_j}|^2\bigl(p_j\cos p_j-q_l\cos p_j-p_j\cos q_l+q_l\cos q_l\bigr)=-\mathfrak D .
\]
The cosine decreases on $(0,\pi)$, and
$(p-q)(\cos q-\cos p)\geq\min(\sin\psi_-,\sin\psi_+)(p-q)^2$ for
$p,q\in[\psi_-,\psi_+]$. The part with $\delta$ consists of four traces for
each arrow, each at most $\norm{\beta_a}^2\norm\delta$.
\end{proof}

For Kronecker quivers the monotonicity of the mass on representations is
\cite[Proposition~9.1]{wue26b}.

\begin{proposition}[Mass of a metric and mass of the filtration]\label{prop:hnmass}
For every tuple of metrics on $\mathcal H$,
\[
 \mass(\mathfrak h)\geq\mass_{\mathrm{HN}}(\mathcal H)
 +\sum_{k=1}^{\mathsf n-1}\bigl(\cos\phi_{k+1}-\cos\phi_k\bigr)\,\mathcal Y_k ,
\]
and the coefficients $\cos\phi_{k+1}-\cos\phi_k$ are positive.
\end{proposition}

\begin{proof}
Let $Q_k=\Pi_k-\Pi_{k-1}$. By the pinching inequality for the trace norm and
$\tr|X|\geq|\tr X|$,
$\mass\geq\sum_k\bigl|\sum_i\tr(Q_{k,i}(z_i+\mu_i))\bigr|$. For a
subrepresentation with projections $\Pi$ one has
$\sum_i\tr(\Pi_i\mu_i)=\sum_a(\norm{\beta_a\Pi_{\mathsf s(a)}}^2-\norm{\Pi_{\mathsf t(a)}\beta_a}^2)
=-\sum_a\norm{\Pi_{\mathsf t(a)}\beta_a(1-\Pi_{\mathsf s(a)})}^2$, as in the
proof of Lemma~\ref{lem:king}(b). So with
$\mathcal Y_0=\mathcal Y_{\mathsf n}=0$,
$\sum_i\tr(Q_{k,i}(z_i+\mu_i))=Z(\operatorname{gr}^k)-(\mathcal Y_k-\mathcal Y_{k-1})$,
and the modulus of this number is at least its component in the direction of
$Z(\operatorname{gr}^k)$, which is
$|Z(\operatorname{gr}^k)|-(\mathcal Y_k-\mathcal Y_{k-1})\cos\phi_k$. Sum
over $k$ and sum by parts.
\end{proof}

\subsection{Unstable cohomology}
Assume (P). The arrows $\beta(t)=(\beta_a(t))$ in the frames of
\eqref{eq:betaflow} form a bounded curve. Let $\omega$ be the set of its
limit points as $t\to\infty$; it is nonempty, compact and connected.

\begin{lemma}\label{lem:limitset}
Assume \emph{(P)}.
\begin{enumerate}
\item[(a)] The mass $\mass(\mathfrak h(t))$ converges, and
 $\beta_a\Phi^{\mathcal H}_{\mathsf s(a)}-\Phi^{\mathcal H}_{\mathsf t(a)}\beta_a$,
 $\dot\beta_a$ and the derivatives of the $\Phi^{\mathcal H}_i$ tend to
 zero.
\item[(b)] There are numbers $\pi>\phi_1'>\dots>\phi_{\mathsf n'}'>0$ and
 nonzero dimension vectors $d_k=(n_{i,k})_i$ with $\sum_kd_k=(n_i)$ and
 $\arg Z(d_k)=\phi_k'$ such that for every $\beta^\infty\in\omega$ the phase
 of $\beta^\infty$ at the vertex $i$ has the eigenvalues $\phi_k'$ with
 multiplicities $n_{i,k}$, and the representation with the arrows
 $\beta^\infty$ is the orthogonal direct sum of representations
 $\mathsf P_k$ of dimension vectors $d_k$, supported on the eigenspaces,
 with the moment map of a harmonic metric of phase $\phi_k'$ on
 $\mathsf P_k$. In particular $\mathsf P_k$ is polystable of phase
 $\phi_k'$.
\item[(c)] The eigenvalues of $\Phi^{\mathcal H}_i(t)$, with multiplicity,
 converge to the $\phi_k'$ with multiplicities $n_{i,k}$, and
 $\mass(\mathfrak h(t))\to\sum_k|Z(d_k)|$.
\end{enumerate}
\end{lemma}

\begin{proof}
(a) By Proposition~\ref{prop:dissipation},
$\mass+4\bar\cN\int_t^\infty\norm\delta$ is nonincreasing and nonnegative,
so $\mass$ converges and $\mathfrak D$ is integrable. By
\eqref{eq:betaflow} the $\dot\beta_a$ are bounded, hence so are the
derivatives of the $\mu_i$ and, by the Lipschitz property, of the
$\Phi^{\mathcal H}_i$. So
$\sum_a\norm{\beta_a\Phi^{\mathcal H}_{\mathsf s(a)}-\Phi^{\mathcal H}_{\mathsf t(a)}\beta_a}^2$
is Lipschitz and integrable and tends to zero; then $\dot\beta_a\to0$ by
\eqref{eq:betaflow}, and the derivatives of $\mu_i$ and
$\Phi^{\mathcal H}_i$ tend to zero.

(b) For $\beta^\infty\in\omega$, by (a) and continuity each arrow
intertwines the phases at its source and target. So the eigenspaces of the
phase for a fixed eigenvalue $p$ form an orthogonal direct summand
$\mathsf P_p$, and on it $\arg(z_i+\mu_i)=p$ for all $i$: the metric is
harmonic of phase $p$ on $\mathsf P_p$, which is therefore polystable
(Lemma~\ref{lem:king}(b)), and $\arg Z(\mathsf P_p)=p$ because
$\sum_i\tr(\Pi_i\mu_i)=0$ for an orthogonal direct summand. Thus the
eigenvalues of the phase of $\beta^\infty$ lie in the finite set of
arguments of the numbers $Z(d)$, $0\neq d\leq(n_i)$. The ordered list of
these eigenvalues at each vertex depends continuously on $\beta$, takes
values in a finite set on $\omega$, and $\omega$ is connected; so it is
constant on $\omega$.

(c) The distance from $\beta(t)$ to $\omega$ tends to zero, and on
$\mathsf P_k$ all eigenvalues of $z_i+\mu_i$ lie on one ray.
\end{proof}

By (c) there is a time after which the spectrum of each
$\Phi^{\mathcal H}_i(t)$ consists of separated clusters near the
$\phi_k'$. Let $P_k(t)$ be the spectral projection of the $k$-th cluster and
$P_{\leq k}=P_1+\dots+P_k$. Since the derivatives of the
$\Phi^{\mathcal H}_i$ tend to zero and the clusters stay separated,
\begin{equation}\label{eq:slow}
 \norm{\dot P_{\leq k}(t)}\to0,\qquad\norm{P_l\beta_aP_k}\to0\quad(l\neq k),
\end{equation}
the second because $X=P_l\beta_aP_k$ satisfies
$X\Phi^{\mathcal H}_{\mathsf s(a)}-\Phi^{\mathcal H}_{\mathsf t(a)}X=P_l(\beta_a\Phi^{\mathcal H}_{\mathsf s(a)}-\Phi^{\mathcal H}_{\mathsf t(a)}\beta_a)P_k$
with separated spectra on the two sides.

\begin{lemma}\label{lem:dichotomy}
Let $\Psi(t)$, $t\geq T$, be a continuous family of Hermitian
$n\times n$ matrices and $P(t)$ a continuously differentiable family of
orthogonal projections of constant rank. Suppose there are numbers
$a_+>a_-$ and $\epsilon\geq0$ with $12\epsilon\leq a_+-a_-$ such that for
all $t\geq T$ and all vectors $w$, with $w'=P(t)w$ and $w''=w-w'$,
\[
 \Rea\inner{w'}{\Psi w}\geq a_+|w'|^2-\epsilon|w'||w''|,\qquad
 \Rea\inner{w''}{\Psi w}\leq a_-|w''|^2+\epsilon|w'||w''|,\qquad
 \norm{\dot P}\leq\epsilon .
\]
Then $\CC^n=S^{\mathrm f}\oplus S^{\mathrm s}$ with
$\dim S^{\mathrm f}=\operatorname{rank}P$, such that for the solutions of
$\dot w=-\frac12\Psi w$ on $[T,\infty)$:
\begin{enumerate}
\item[(i)] if $w(T)\in S^{\mathrm f}$, then
 $|w''(t)|\leq\frac{6\epsilon}{a_+-a_-}|w'(t)|$ and
 $|w(t)|^2\leq2e^{-(a_+-3\epsilon)(t-T)}|w(T)|^2$ for $t\geq T$;
\item[(ii)] if $w(T)\in S^{\mathrm s}$, then
 $|w(t)|^2\geq\frac12e^{-(a_-+3\epsilon)(t-T)}|w(T)|^2$ for $t\geq T$.
\end{enumerate}
\end{lemma}

\begin{proof}
Since $P\dot PP=0=(1-P)\dot P(1-P)$, a solution satisfies
$\frac d{dt}|w'|^2\leq-a_+|w'|^2+3\epsilon|w'||w''|$ and
$\frac d{dt}|w''|^2\geq-a_-|w''|^2-3\epsilon|w'||w''|$. Let
$c_0=6\epsilon/(a_+-a_-)\leq\frac12$. For $c\in[c_0,1]$ with $c>0$ and a
nonzero solution, at a time with $|w''|=c|w'|$,
\[
 \frac d{dt}\bigl(c^2|w'|^2-|w''|^2\bigr)\leq c^2|w'|^2\bigl(-(a_+-a_-)+3\epsilon(c+c^{-1})\bigr)<0,
\]
because $3\epsilon c\leq\frac14(a_+-a_-)$ and
$3\epsilon c^{-1}\leq\frac12(a_+-a_-)$. So a solution in the open cone
$\{|w''|>c|w'|\}$ cannot reach its boundary: these cones are invariant
forwards in time. Let $S^{\mathrm s}=\ker P(T)$. A solution starting there
stays in the cone with $c=1$, where
$\frac d{dt}|w''|^2\geq-(a_-+3\epsilon)|w''|^2$; this gives (ii). For
$\tau\geq T$ let $S_\tau$ be the space of the initial values $w(T)$ of the
solutions with $w(\tau)$ in the range of $P(\tau)$. Such a solution
satisfies $|w''|\leq c_0|w'|$ on $[T,\tau]$, since otherwise it would lie, at
the time $\tau$, in one of the open cones with $c\in[c_0,1]$, $c>0$. Let $S^{\mathrm f}$ be a
limit point of the spaces $S_\tau$ in the Grassmannian as $\tau\to\infty$.
By continuity every solution starting in $S^{\mathrm f}$ satisfies
$|w''|\leq c_0|w'|$ for all $t\geq T$, and then
$\frac d{dt}|w'|^2\leq-(a_+-3\epsilon)|w'|^2$ and $|w|^2\leq2|w'|^2$; this is
(i). The two growth estimates exclude a common nonzero vector, and the
dimensions add up to $n$.
\end{proof}

\begin{theorem}[Harder--Narasimhan asymptotics]\label{thm:hn}
Let $\mathcal H$ be a nonzero representation of $Q$, with
Harder--Narasimhan filtration $(\mathcal H^k)$, factors
$\operatorname{gr}^k$ and phases $\phi_1>\dots>\phi_{\mathsf n}$, and let
$\mathfrak h(t)$ satisfy \emph{(P)}. Then:
\begin{enumerate}
\item[(i)] the eigenvalues of $\Phi^{\mathcal H}_i(t)$, with multiplicity,
 converge to the $\phi_k$ with multiplicities
 $\dim\operatorname{gr}^k_i$, and $\mass(\mathfrak h(t))$ converges to
 $\mass_{\mathrm{HN}}(\mathcal H)$, the infimum of the mass over all metrics
 on $\mathcal H$;
\item[(ii)] for every nonzero $v\in\mathcal H_i$,
 \[
 \lim_{t\to\infty}\frac1t\log\mathfrak h_i(t)(v,v)=\varphi-\phi_k,
 \]
 where $k$ is the least index with $v\in\mathcal H^k_i$; so the
 Harder--Narasimhan filtration is the filtration of the spaces by the
 exponential rate of decay of the metric;
\item[(iii)] for every $\epsilon>0$ there is a constant $C$ with
 $\mathcal Y_k(t)\leq Ce^{-2\varkappa_kt}$,
 $\varkappa_k=\min\{\frac12(\phi_k-\phi_{k+1})-\epsilon,\ \kappa\}$, for
 $k<\mathsf n$ (with $\kappa=\infty$ if $\delta=0$);
\item[(iv)] for each $k$ the metric induced by $\mathfrak h(t)$ on
 $\operatorname{gr}^k$, multiplied by $e^{(\phi_k-\varphi)t}$, satisfies
 \emph{(P)} for the semistable representation $\operatorname{gr}^k$. For one
 factor the original perturbation and its rate $\kappa$ are retained.
 For two or more factors take $0<2\epsilon<\min_{l<\mathsf n}(\phi_l-\phi_{l+1})$;
 the perturbation then has rate $\min_{l<\mathsf n}\varkappa_l$. Hence
 Proposition~\ref{prop:perturbed}, Theorem~\ref{thm:comparison} and
 Corollary~\ref{cor:iterated} apply to it: its moment map tends to the
 stability parameter of $\operatorname{gr}^k$, it is an asymptotic solution
 of the weighted gradient flow of $\operatorname{gr}^k$, and if
 $\operatorname{gr}^k$ is polystable it converges exponentially to a
 harmonic metric of $\operatorname{gr}^k$.
\end{enumerate}
\end{theorem}

\begin{proof}
\emph{Step 1: the filtration by exponents.} Fix $k<\mathsf n'$, a vertex
$i$, and $\epsilon>0$ with $60\epsilon\leq\phi_k'-\phi_{k+1}'$. By
Lemma~\ref{lem:limitset} and \eqref{eq:slow} there is a time $T$ after which
the clusters are within $\epsilon$ of the $\phi_l'$,
$\norm\delta\leq\epsilon$ and $\norm{\dot P_{\leq k}}\leq\epsilon$. For a
vector $v\in\mathcal H_i$ the coordinates $w$ of $v$ in the frame satisfy
$|w|^2=\mathfrak h_i(v,v)$ and $\dot w=-\frac12\Psi w$ with
$\Psi=\Phi^{\mathcal H}_i-\varphi+\delta_i$. Since $P=P_{\leq k}$ commutes
with $\Phi^{\mathcal H}_i$, the hypotheses of Lemma~\ref{lem:dichotomy} hold
with $a_+=\phi_k'-\varphi-2\epsilon$ and
$a_-=\phi_{k+1}'-\varphi+2\epsilon$. In terms of vectors of $\mathcal H_i$
the lemma gives complementary subspaces, of which the first has the
dimension $\sum_{l\leq k}n_{i,l}$, with
$\mathfrak h_i(t)(v,v)\leq C_ve^{-(\phi_k'-\varphi-5\epsilon)t}$ on the
first and
$\mathfrak h_i(t)(v,v)\geq c_ve^{-(\phi_{k+1}'-\varphi+5\epsilon)t}$,
$c_v>0$, for the nonzero vectors of the second. A vector with a nonzero
component in the second subspace then satisfies
$\liminf\frac1t\log\mathfrak h_i(t)(v,v)\geq-(\phi_{k+1}'-\varphi+5\epsilon)$.
So the first subspace is the set of all $v$ with
$\limsup\frac1t\log\mathfrak h_i(t)(v,v)\leq-(\phi_k'-\varphi-5\epsilon)$;
these sets decrease with $\epsilon$ and have the same dimension, so they
coincide with
\[
 F_{k,i}:=\Bigl\{v\in\mathcal H_i:\ \limsup_{t\to\infty}\frac1t\log\mathfrak h_i(t)(v,v)\leq\varphi-\phi_k'\Bigr\},
 \qquad\dim F_{k,i}=\sum_{l\leq k}n_{i,l},
\]
and for $v\in F_{k,i}\setminus F_{k-1,i}$ the limit of
$\frac1t\log\mathfrak h_i(t)(v,v)$ exists and equals $\varphi-\phi_k'$
(with $F_0=0$ and $F_{\mathsf n'}=\mathcal H$). Here the lower bound for
$k=1$ and the upper bound for $k=\mathsf n'$, which do not come from a gap,
follow from $\frac d{dt}|w|^2=-\inner w{\Psi w}$ and
Lemma~\ref{lem:limitset}(c): the eigenvalues of $\Psi$ lie eventually between
$\phi_{\mathsf n'}'-\varphi-2\epsilon$ and $\phi_1'-\varphi+2\epsilon$.

\emph{Step 2: the $F_k$ are subrepresentations.} For an arrow $a$ of
$\mathcal H$ and $v\in\mathcal H_{\mathsf s(a)}$, the norm of the image of
$v$ for $\mathfrak h_{\mathsf t(a)}$ is at most $\bar\cN^{1/2}$ times the
norm of $v$ for $\mathfrak h_{\mathsf s(a)}$. So the arrows do not decrease
the rate of decay, and they map $F_k$ to $F_k$. The factor $F_k/F_{k-1}$ has
the dimension vector $d_k$, hence $\arg Z(F_k/F_{k-1})=\phi_k'$.

\emph{Step 3.} By Lemma~\ref{lem:dichotomy}(i), applied with $\epsilon$
arbitrarily small for large $T$, the coordinates of the vectors of $F_k$
form the graph of a linear map from the range of $P_{\leq k}(t)$ to its
orthogonal complement whose norm tends to zero. So for the
$\mathfrak h(t)$-orthogonal projections $\Pi_k'(t)$ onto $F_k$, in the
frames, $\norm{\Pi_k'-P_{\leq k}}\to0$; and with
$Q_k=\Pi_k'-\Pi_{k-1}'$, by \eqref{eq:slow} and Lemma~\ref{lem:limitset}(c),
\begin{equation}\label{eq:blockslimit}
 Q_l\beta_aQ_k\to0\ (l\neq k),\qquad
 \norm{Q_k(\Phi^{\mathcal H}-\phi_k')Q_k}\to0,\qquad
 Q_l\Phi^{\mathcal H}Q_k\to0\ (l\neq k).
\end{equation}

\emph{Step 4: the factors are semistable.} In frames adapted to the
decomposition $1=\sum_kQ_k(t)$ the arrows are block upper triangular, and the
$k$-th diagonal block is the tuple of arrows of the representation
$F_k/F_{k-1}$ in orthonormal frames of its induced metric, a point of the
orbit of this representation under the changes of bases. Along a suitable
sequence of times these frames and the arrows converge, and by
\eqref{eq:blockslimit} the $k$-th diagonal block converges to the arrows of
the summand $\mathsf P_k$ of a point of $\omega$. So the closure of the orbit
of $F_k/F_{k-1}$ contains a representation isomorphic to $\mathsf P_k$, which
is semistable of phase $\phi_k'$. If $F_k/F_{k-1}$ had a subrepresentation
of a dimension vector $d'$ with $\arg Z(d')>\phi_k'$, then so would every
representation in the closure of its orbit, because the set of
representations with a subrepresentation of dimension vector $d'$ is the
image of a closed set under a projection with compact fibres, a product of
Grassmannians. This contradicts the semistability of $\mathsf P_k$. So
$F_k/F_{k-1}$ is semistable of phase $\phi_k'$.

\emph{Step 5.} The filtration $(F_k)$ has semistable factors of strictly
decreasing phases. Hence it is the Harder--Narasimhan filtration:
$\mathsf n'=\mathsf n$, $F_k=\mathcal H^k$, $\phi_k'=\phi_k$ and
$d_k=\dim\operatorname{gr}^k$. With Lemma~\ref{lem:limitset}(c) and Step~1
this gives (i) and (ii); the mass of the filtration is the infimum of the
mass by Proposition~\ref{prop:hnmass}.

\emph{Step 6: the second fundamental forms.} Fix $k<\mathsf n$ and write
$\Pi=\Pi_k$. For an arrow $a$ let
$\upsilon_a=\Pi_{\mathsf t(a)}\beta_a(1-\Pi_{\mathsf s(a)})$, and for a
Hermitian matrix $M$ at a vertex let $M^{\mathrm{off}}=\Pi M(1-\Pi)+(1-\Pi)M\Pi$.
With $\tilde\Phi=\Phi^{\mathcal H}+\delta$ we claim that
\begin{equation}\label{eq:secondform}
 \frac{d\mathcal Y_k}{dt}=\sum_a\Bigl[\tr\bigl(\upsilon_a^\dagger\upsilon_a\,\tilde\Phi_{\mathsf s(a)}\bigr)
 -\tr\bigl(\upsilon_a\upsilon_a^\dagger\,\tilde\Phi_{\mathsf t(a)}\bigr)\Bigr]
 +\sum_i\tr\bigl(\tilde\Phi_i^{\mathrm{off}}\mu_i^{\mathrm{off}}\bigr).
\end{equation}
Indeed, the coordinate maps of the frames satisfy
$\dot gg^{-1}=-\frac12(\tilde\Phi-\varphi)$, a Hermitian matrix, and the
range of $\Pi$ is the image of a fixed subspace, so
$\dot\Pi=-\frac12\tilde\Phi^{\mathrm{off}}$. Differentiate
$\upsilon_a$ with \eqref{eq:betaflow}, use
$(1-\Pi_{\mathsf t(a)})\beta_a\Pi_{\mathsf s(a)}=0$, and take the trace
against $\upsilon_a^\dagger$: with
$\beta_a'=\Pi_{\mathsf t(a)}\beta_a\Pi_{\mathsf s(a)}$ and
$\beta_a''=(1-\Pi_{\mathsf t(a)})\beta_a(1-\Pi_{\mathsf s(a)})$,
\begin{align*}
 \frac d{dt}\norm{\upsilon_a}^2={}&\tr\bigl(\upsilon_a^\dagger\upsilon_a\tilde\Phi_{\mathsf s(a)}\bigr)
 -\tr\bigl(\upsilon_a\upsilon_a^\dagger\tilde\Phi_{\mathsf t(a)}\bigr)\\
 &-2\Rea\tr\bigl(\Pi\tilde\Phi_{\mathsf t(a)}(1-\Pi)\beta_a''\upsilon_a^\dagger\bigr)
 +2\Rea\tr\bigl(\Pi\tilde\Phi_{\mathsf s(a)}(1-\Pi)\upsilon_a^\dagger\beta_a'\bigr),
\end{align*}
and
$(1-\Pi)\mu_i\Pi=\sum_{\mathsf s(a)=i}\upsilon_a^\dagger\beta_a'-\sum_{\mathsf t(a)=i}\beta_a''\upsilon_a^\dagger$;
summing over the arrows gives \eqref{eq:secondform}.

The term $\sum_i\tr((\Phi^{\mathcal H}_i)^{\mathrm{off}}\mu_i^{\mathrm{off}})$
is nonpositive: the matrix $2\Pi-1$ is unitary and Hermitian, conjugation
by it changes the sign of the off-diagonal parts, and
$\Phi^{\mathcal H}_i$ is a decreasing function of $\mu_i$, so four times the
trace is $\tr((g(A)-g(B))(A-B))\leq0$ with $A=\mu_i$,
$B=(2\Pi-1)\mu_i(2\Pi-1)$ and $g=\arg(z_i+\,\cdot\,)$. By
\eqref{eq:blockslimit} and Step~5, for every $\epsilon>0$ and large $t$
the eigenvalues of $\Pi\Phi^{\mathcal H}\Pi$ on the range of $\Pi$ are at
least $\phi_k-\epsilon$ and those of $(1-\Pi)\Phi^{\mathcal H}(1-\Pi)$ on the
complement are at most $\phi_{k+1}+\epsilon$. The terms with $\delta$ are at
most $2\norm\delta\mathcal Y_k+2\sqrt2\,\norm\delta(\bar\cN\mathcal Y_k)^{1/2}$.
Hence
\[
 \frac{d\mathcal Y_k}{dt}\leq-\bigl(\phi_k-\phi_{k+1}-2\epsilon-2\norm\delta\bigr)\mathcal Y_k
 +2\sqrt2\,\norm\delta\,(\bar\cN\mathcal Y_k)^{1/2}
\]
for large $t$. With $\norm\delta\leq c_\delta e^{-\kappa(t-\tau_0)}$ this
gives (iii); one applies the inequality with $\epsilon/2$ in the place of
$\epsilon$, so that a factor $t$ arising when the two exponents coincide is
absorbed.

\emph{Step 7: the factors.} If $\mathsf n=1$, the factor is $\mathcal H$
and $\phi_1=\varphi$, so (iv) follows directly from (P). For $\mathsf n>1$,
the restriction of a metric to a subspace and
the quotient metric move with the compressed velocity. So in the adapted
frames of Step~4 the metric induced on $\operatorname{gr}^k$, multiplied by
$e^{(\phi_k-\varphi)t}$, has the velocity
$\phi_k-Q_k\Phi^{\mathcal H}Q_k-Q_k\delta Q_k$, and its arrows are the $k$-th
diagonal block. Let $\mu^{(k)}$ be the moment map of this block, $D$ the
block-diagonal part of $\mu$ and $O=\mu-D$. Every block of $O$, and every
block of $Q_kDQ_k-\mu^{(k)}$, is a sum of products of two blocks of the
arrows of which at least one, respectively both, are off-diagonal; so with
$Y^2=\sum_a\sum_{l<k}\norm{Q_l\beta_aQ_k}^2\leq\sum_k\mathcal Y_k$,
$\norm O\leq4\bar\cN^{1/2}Y$ and $\norm{Q_kDQ_k-\mu^{(k)}}\leq Y^2$. By the
Lipschitz property, and because a function of $D$ is block diagonal,
\[
 \norm{Q_k\Phi^{\mathcal H}Q_k-\arg(z+\mu^{(k)})}\leq\frac{4\bar\cN^{1/2}Y+Y^2}{\min_iy_i}.
\]
So the velocity is $\phi_k-\arg(z+\mu^{(k)})-\delta^{(k)}$ with
$\delta^{(k)}$ exponentially small by (iii), and the squared norms of the
arrows of the block sum to at most $\bar\cN$: this is (P) for
$\operatorname{gr}^k$, which is semistable of phase $\phi_k$.
\end{proof}

\begin{remark}\label{rem:hn}
(1) For a subrepresentation of a representation, and $\delta=0$, the
identity \eqref{eq:secondform} holds at every time, with $\mathcal Y$ the
squared norm of its second fundamental form. Its last term is nonpositive.
So the second fundamental form decays as long as the phases on the
subrepresentation exceed those on its orthogonal complement.

(2) The filtration by the rates of decay is a filtration by
subrepresentations for the simple reason of Step~2; the theorem identifies
the rates and the dimensions. For the gradient flow the limiting filtration
is described in the same terms in \cite[\S1.1]{hkkp23}, and for unstable
points of Kempf--Ness gradient flows the term linear in $t$ is in
\cite{lam26}.

(3) By (i) and Proposition~\ref{prop:dissipation}, along the phase flow of a
representation of a quiver without oriented cycles the mass decreases to the
mass of the object for the stability condition; the largest and the
smallest eigenvalue of the phase tend to $\phi_1$ and $\phi_{\mathsf n}$.
\end{remark}
